% ============================================================
% Activité de découverte — Ch05 Équations du premier degré
% 2nde Bac Pro MAMA — Mathématiques — Durée : 35 min
% Compilable : pdflatex (Overleaf)
% ============================================================
\documentclass[a4paper,11pt]{article}

% --- Encodage et langue ---
\usepackage[utf8]{inputenc}
\usepackage[T1]{fontenc}
\usepackage[french]{babel}

% --- Mise en page ---
\usepackage[margin=1.8cm,top=1.5cm,bottom=1.5cm]{geometry}
\setlength{\parindent}{0pt}
\setlength{\parskip}{4pt}

% --- Couleurs et boîtes ---
\usepackage[most]{tcolorbox}
\usepackage{xcolor}

% Palette (maths/seconde)
\definecolor{prim}{HTML}{0056b3}
\definecolor{primbg}{HTML}{ebf5ff}
\definecolor{primborder}{HTML}{bee3f8}
\definecolor{greenbox}{HTML}{dcfce7}
\definecolor{greenborder}{HTML}{16a34a}
\definecolor{redborder}{HTML}{ef4444}
\definecolor{orangebg}{HTML}{fff7ed}
\definecolor{orangeborder}{HTML}{f97316}
\definecolor{purplebg}{HTML}{f5f0ff}
\definecolor{purpleborder}{HTML}{7c3aed}
\definecolor{yellowbg}{HTML}{fefce8}
\definecolor{yellowborder}{HTML}{ca8a04}
\definecolor{graybg}{HTML}{f3f4f6}
\definecolor{grayborder}{HTML}{9ca3af}
\definecolor{situationbg}{HTML}{faf5ff}
\definecolor{situationborder}{HTML}{a78bfa}
\definecolor{retenirbg}{HTML}{f0fdf4}
\definecolor{retenirborder}{HTML}{22c55e}

% --- Mathématiques ---
\usepackage{amsmath,amssymb}

% --- Tableaux ---
\usepackage{array,booktabs}

% --- Listes ---
\usepackage{enumitem}
\setlist{nosep}

% --- Divers ---
\usepackage{graphicx}

% ============================================================
% Badges de compétences
% ============================================================
\newtcbox{\badgeAPP}{
  on line,colback=primbg,colframe=prim,
  boxrule=0.6pt,arc=3pt,
  left=3pt,right=3pt,top=1pt,bottom=1pt,
  fontupper=\scriptsize\bfseries\color{prim}
}
\newtcbox{\badgeANA}{
  on line,colback=purplebg,colframe=purpleborder,
  boxrule=0.6pt,arc=3pt,
  left=3pt,right=3pt,top=1pt,bottom=1pt,
  fontupper=\scriptsize\bfseries\color{purpleborder}
}
\newtcbox{\badgeREA}{
  on line,colback=greenbox,colframe=greenborder,
  boxrule=0.6pt,arc=3pt,
  left=3pt,right=3pt,top=1pt,bottom=1pt,
  fontupper=\scriptsize\bfseries\color{greenborder}
}
\newtcbox{\badgeVAL}{
  on line,colback=orangebg,colframe=orangeborder,
  boxrule=0.6pt,arc=3pt,
  left=3pt,right=3pt,top=1pt,bottom=1pt,
  fontupper=\scriptsize\bfseries\color{orangeborder}
}
\newtcbox{\badgeCOM}{
  on line,colback=yellowbg,colframe=yellowborder,
  boxrule=0.6pt,arc=3pt,
  left=3pt,right=3pt,top=1pt,bottom=1pt,
  fontupper=\scriptsize\bfseries\color{yellowborder}
}

% ============================================================
% Boîtes tcolorbox
% ============================================================

% Situation professionnelle
\newtcolorbox{situation}{
  colback=situationbg,colframe=situationborder,
  fonttitle=\bfseries,title=Situation professionnelle — Atelier de menuiserie,
  boxrule=1pt,arc=4pt,
  left=8pt,right=8pt,top=6pt,bottom=6pt,
  toptitle=4pt,bottomtitle=4pt,
  borderline west={3pt}{0pt}{situationborder}
}

% Problématique
\newtcolorbox{problematique}{
  colback=primbg,colframe=prim,
  fonttitle=\bfseries,title=Problématique,
  boxrule=1pt,arc=4pt,
  left=8pt,right=8pt,top=4pt,bottom=4pt,
  toptitle=3pt,bottomtitle=3pt,
  borderline west={4pt}{0pt}{prim}
}

% Question
\newtcolorbox{question}[1][]{
  colback=white,colframe=prim,
  boxrule=0.8pt,arc=3pt,
  left=8pt,right=8pt,top=5pt,bottom=5pt,
  toptitle=3pt,bottomtitle=3pt,
  fonttitle=\bfseries,
  title={#1}
}

% Correction
\newtcolorbox{correction}{
  colback=graybg,colframe=grayborder,
  boxrule=0.5pt,arc=3pt,
  left=8pt,right=8pt,top=4pt,bottom=4pt,
  fontupper=\small
}

% À retenir
\newtcolorbox{retenir}{
  colback=retenirbg,colframe=retenirborder,
  fonttitle=\bfseries\large,title=À retenir,
  boxrule=1.2pt,arc=4pt,
  left=8pt,right=8pt,top=6pt,bottom=6pt,
  toptitle=4pt,bottomtitle=4pt,
  borderline west={4pt}{0pt}{retenirborder}
}

% Objectifs
\newtcolorbox{objectifs}{
  colback=primbg,colframe=primborder,
  boxrule=0.8pt,arc=3pt,
  left=8pt,right=8pt,top=4pt,bottom=4pt
}

% ============================================================
\begin{document}

% --- EN-TÊTE ---
{\centering
  {\LARGE\bfseries\color{prim} Activité de découverte — Équations du premier degré}\\[6pt]
  {\normalsize Chapitre 5 \;$\vert$\; Seconde Bac Pro \;$\vert$\; Mathématiques \;$\vert$\; \raisebox{-0.5pt}{\large$\circlearrowleft$}\;35 min}\\[10pt]
}

% --- OBJECTIFS ---
\begin{objectifs}
\textbf{Objectifs :}
\begin{itemize}[leftmargin=14pt]
  \item Découvrir la notion d'équation du premier degré à partir d'une situation concrète
  \item Apprendre à traduire un problème professionnel en équation
  \item Construire la méthode de résolution étape par étape
\end{itemize}
\end{objectifs}

\vspace{6pt}

% ============================================================
% SITUATION PROFESSIONNELLE
% ============================================================
\begin{situation}
\textbf{Théo Rimbaud} est artisan menuisier dans l'entreprise \textbf{Bois \& Agencement\,33}, à Bordeaux.

\medskip
Il dispose d'une \textbf{planche de chêne brut} de \textbf{240\,cm} de longueur.
Il doit y découper les pièces suivantes :

\medskip
\begin{center}
\renewcommand{\arraystretch}{1.3}
\begin{tabular}{|l|c|c|}
\hline
\textbf{Pièce} & \textbf{Quantité} & \textbf{Longueur}\\
\hline
Étagères identiques & 4 & \textbf{? cm}\\
\hline
Tasseau de renfort & 1 & 30 cm\\
\hline
Chute (incompressible) & — & 18 cm\\
\hline
\textbf{Total} & — & \textbf{240 cm}\\
\hline
\end{tabular}
\end{center}
\end{situation}

\vspace{4pt}

% --- PROBLÉMATIQUE ---
\begin{problematique}
\textit{Quelle longueur doit avoir chaque étagère pour utiliser toute la planche~?}
\end{problematique}

\vspace{8pt}

% ============================================================
% QUESTIONS
% ============================================================

% ---- Q1 ----
\begin{question}[Question 1 \quad \badgeAPP{APP}]
En lisant la situation, identifier :
\begin{enumerate}[label=\alph*)]
  \item la longueur totale de la planche~;
  \item le nombre d'étagères à découper~;
  \item les longueurs connues et la longueur inconnue.
\end{enumerate}
\end{question}

\begin{correction}
\textbf{a)} La planche mesure \textbf{240\,cm} au total.\\
\textbf{b)} Il faut découper \textbf{4 étagères} identiques.\\
\textbf{c)} Longueurs connues : tasseau = 30\,cm, chute = 18\,cm. Longueur inconnue : la longueur d'une étagère.
\end{correction}

\vspace{6pt}

% ---- Q2 ----
\begin{question}[Question 2 \quad \badgeAPP{APP}]
On note $x$ la longueur (en cm) d'une étagère.
\begin{enumerate}[label=\alph*)]
  \item Pourquoi $x$ est-elle une \textbf{inconnue}~?
  \item Exprimer en fonction de $x$ la longueur totale occupée par les 4~étagères.
\end{enumerate}
\end{question}

\begin{correction}
\textbf{a)} $x$ est une inconnue car on ne connaît pas (encore) sa valeur : c'est la grandeur à déterminer.\\
\textbf{b)} Les 4 étagères mesurent chacune $x$\,cm, donc la longueur totale des étagères est $4 \times x = 4x$\;cm.
\end{correction}

\vspace{6pt}

% ---- Q3 ----
\begin{question}[Question 3 \quad \badgeANA{ANA}]
En utilisant le fait que la somme de toutes les pièces donne la longueur totale de la planche, écrire une \textbf{égalité} reliant $x$ aux longueurs connues.
\end{question}

\begin{correction}
La planche est découpée en : 4 étagères ($4x$) + 1 tasseau (30\,cm) + 1 chute (18\,cm) = 240\,cm.

\medskip
L'égalité obtenue est :
\[ 4x + 30 + 18 = 240 \]

C'est une \textbf{équation du premier degré} d'inconnue $x$.
\end{correction}

\vspace{6pt}

% ---- Q4 ----
\begin{question}[Question 4 \quad \badgeANA{ANA}]
Simplifier le membre de gauche en regroupant les termes \textbf{sans $x$} (les nombres connus).
\end{question}

\begin{correction}
On additionne les constantes : $30 + 18 = 48$.

L'équation devient :
\[ 4x + 48 = 240 \]
\end{correction}

\vspace{6pt}

% ---- Q5 ----
\begin{question}[Question 5 \quad \badgeREA{REA}]
\begin{enumerate}[label=\alph*)]
  \item Pour isoler le terme $4x$, soustraire 48 des \textbf{deux côtés} de l'équation. Écrire le résultat.
  \item Pour trouver $x$, diviser les deux côtés par 4. Quelle est la valeur de $x$~?
\end{enumerate}
\end{question}

\begin{correction}
\textbf{a)} On soustrait 48 de chaque côté :
\[ 4x + 48 - 48 = 240 - 48 \]
\[ 4x = 192 \]

\textbf{b)} On divise chaque côté par 4 :
\[ x = \frac{192}{4} = 48 \]

Chaque étagère mesure \textbf{48\,cm}.
\end{correction}

\vspace{6pt}

% ---- Q6 ----
\begin{question}[Question 6 \quad \badgeVAL{VAL}]
\begin{enumerate}[label=\alph*)]
  \item Vérifier en remplaçant $x$ par 48 dans l'équation de départ : calculer $4 \times 48 + 30 + 18$.
  \item Le résultat obtenu est-il bien égal à 240~?
  \item Rédiger une phrase de conclusion répondant à la problématique.
\end{enumerate}
\end{question}

\begin{correction}
\textbf{a)} $4 \times 48 + 30 + 18 = 192 + 30 + 18 = 240$.\\
\textbf{b)} Oui, on obtient bien 240. La solution est donc correcte.\\
\textbf{c)} Chaque étagère doit mesurer \textbf{48\,cm} pour utiliser toute la planche de 240\,cm.
\end{correction}

\vspace{6pt}

% ---- Q7 ----
\begin{question}[Question 7 \quad \badgeREA{REA}]
\textbf{Nouveau problème.} Théo dispose maintenant d'une planche de \textbf{3\,m}. Il doit découper \textbf{5 étagères identiques}, un tasseau de \textbf{40\,cm} et il y aura une chute de \textbf{20\,cm}.

\medskip
En utilisant la même méthode, trouver la longueur de chaque étagère.
\end{question}

\begin{correction}
Convertir : 3\,m = 300\,cm.

Inconnue : $x$ = longueur d'une étagère (en cm).

Équation : $5x + 40 + 20 = 300$, soit $5x + 60 = 300$.

Résolution :
\[ 5x = 300 - 60 = 240 \]
\[ x = \frac{240}{5} = 48 \text{\;cm} \]

Vérification : $5 \times 48 + 40 + 20 = 240 + 60 = 300$\;\checkmark

Chaque étagère mesure \textbf{48\,cm}.
\end{correction}

\vspace{6pt}

% ---- Q8 ----
\begin{question}[Question 8 \quad \badgeREA{REA}]
\textbf{Autre contexte.} Théo commande des poignées de tiroir à \textbf{3{,}50\,\texteuro{}} pièce. Les frais de port sont de \textbf{8\,\texteuro{}}. Le total de la facture est de \textbf{57\,\texteuro{}}.

\medskip
Combien de poignées a-t-il commandées~? {\small(Poser l'équation, résoudre, vérifier et conclure.)}
\end{question}

\begin{correction}
Inconnue : $x$ = nombre de poignées commandées.

Équation : $3{,}50\,x + 8 = 57$.

Résolution :
\[ 3{,}50\,x = 57 - 8 = 49 \]
\[ x = \frac{49}{3{,}50} = 14 \]

Vérification : $3{,}50 \times 14 + 8 = 49 + 8 = 57$\;\checkmark

Théo a commandé \textbf{14 poignées}.
\end{correction}

\vspace{6pt}

% ---- Q9 ----
\begin{question}[Question 9 \quad \badgeANA{ANA}]
Parmi les expressions suivantes, lesquelles sont des \textbf{équations du premier degré} en~$x$~? Justifier.

\medskip
\begin{enumerate*}[label=\textbf{\arabic*.}\;,itemjoin=\qquad]
  \item $3x + 5 = 14$
  \item $x^2 = 9$
  \item $2x - 7 = x + 1$
  \item $\dfrac{1}{x} = 4$
  \item $5x = 0$
  \item $x^2 + x = 3$
\end{enumerate*}
\end{question}

\begin{correction}
\textbf{1.} $3x + 5 = 14$ \;\checkmark\; Oui — $x$ au degré 1.\\
\textbf{2.} $x^2 = 9$ \;$\boldsymbol{\times}$\; Non — $x$ au degré 2 (équation du second degré).\\
\textbf{3.} $2x - 7 = x + 1$ \;\checkmark\; Oui — $x$ au degré 1 des deux côtés.\\
\textbf{4.} $\frac{1}{x} = 4$ \;$\boldsymbol{\times}$\; Non — $x$ est au dénominateur, ce n'est pas du degré 1.\\
\textbf{5.} $5x = 0$ \;\checkmark\; Oui — de la forme $ax = b$ avec $b = 0$.\\
\textbf{6.} $x^2 + x = 3$ \;$\boldsymbol{\times}$\; Non — présence de $x^2$ (degré 2).
\end{correction}

\vspace{6pt}

% ---- Q10 ----
\begin{question}[Question 10 \quad \badgeCOM{COM}]
Compléter le texte suivant avec les mots : \textit{premier degré}, \textit{inconnue}, \textit{solution}, \textit{vérifier}, \textit{isoler}.

\medskip
\begin{quote}
\og Une équation du \rule{2.5cm}{0.4pt} est une égalité contenant une \rule{2.5cm}{0.4pt} (souvent notée $x$) qui apparaît au degré 1. Résoudre une équation, c'est chercher la valeur de $x$ — appelée \rule{2.5cm}{0.4pt} — qui rend l'égalité vraie. Pour cela, on effectue des opérations des deux côtés de l'égalité afin d'\rule{2.5cm}{0.4pt} $x$. Il faut ensuite \rule{2.5cm}{0.4pt} le résultat en remplaçant $x$ dans l'équation de départ. \fg
\end{quote}
\end{question}

\begin{correction}
\og Une équation du \textbf{premier degré} est une égalité contenant une \textbf{inconnue} (souvent notée $x$) qui apparaît au degré 1. Résoudre une équation, c'est chercher la valeur de $x$ — appelée \textbf{solution} — qui rend l'égalité vraie. Pour cela, on effectue des opérations des deux côtés de l'égalité afin d'\textbf{isoler} $x$. Il faut ensuite \textbf{vérifier} le résultat en remplaçant $x$ dans l'équation de départ. \fg
\end{correction}

\vspace{10pt}

% ============================================================
% SYNTHÈSE — À RETENIR
% ============================================================
\begin{retenir}
\textbf{Réponse à la problématique :} Chaque étagère doit mesurer \textbf{48\,cm}. On a trouvé cette valeur en traduisant la situation par une équation ($4x + 48 = 240$) puis en isolant $x$.

\bigskip
\textbf{Ce qu'il faut retenir :}
\begin{itemize}[leftmargin=14pt]
  \item Une \textbf{équation du premier degré} est une égalité contenant une inconnue $x$ au degré 1.
  \item \textbf{Résoudre} une équation, c'est trouver la valeur de $x$ (la \textbf{solution}) qui vérifie l'égalité.
  \item Pour résoudre, on effectue la \textbf{même opération des deux côtés} du signe $=$ pour isoler $x$.
  \item On doit toujours \textbf{vérifier} la solution dans l'équation de départ et \textbf{conclure en phrase} dans le contexte du problème.
\end{itemize}

\bigskip
\textbf{Méthode en 5 étapes :}
\begin{enumerate}[leftmargin=18pt]
  \item Choisir l'inconnue $x$ et préciser l'unité.
  \item Traduire le problème en équation.
  \item Résoudre : regrouper les termes, diviser par le coefficient.
  \item Vérifier dans l'équation de départ.
  \item Conclure en phrase avec l'unité.
\end{enumerate}
\end{retenir}

\end{document}
